\chapter[LITERATURE SURVEY]{Literature Survey}

This chapter summarizes the technical decisions made for the project. It starts with credit classification and tree boosting, then moves on to SHAP and LIME and the assessment of local explanations. The discussion concludes with fairness and decision-support research, highlighting the ideas that are supported by the current system and those that need more evidence than the current system offers.

\section{Credit Scoring and Predictive Models}
\subsection{Credit Scoring is a Classification Problem}
Credit scoring isn't just a simple classification job with financial field names. The available records are influenced by previous lending practices, the criteria used to define a successful account and the applicant pool. Hand and Henley discuss statistical methods used in consumer credit scoring and problems that can occur in this context \cite{handHenley1997}.

This background is important to understand because a project with generated approval labels will have to make a narrower claim than a study with observed repayment outcomes. TrustFin AI can analyze the reproduction of the decision rules generated by a classifier. It cannot make that result alone a conclusion that the same classifier would rank actual credit risk correctly.

Lessmann et al. compare the performance of classification algorithms on credit-scoring datasets \cite{lessmann2015}. Their work does not make a general statement that one method is always better than another, but rather it helps to make empirical comparisons. For this project, XGBoost is chosen as the implementation. A comparison with other classifiers is not included in the saved product evaluation and superiority over those other classifiers is not claimed here.

\subsection{Gradient Boosting and XGBoost}
Gradient boosting builds an additive model by a series of learners. Every new learner is assigned to enhance the goal of the current ensemble. The general basis for this approach is given by Friedman's formulation \cite{friedman2001}. For structured loan records, the resulting model can be used to model interactions that would be cumbersome to model as one linear score.

According to Chen and Guestrin, XGBoost is a scalable tree-boosting system, where regularisation and implementation methods are provided for efficient learning \cite{chen2016}. TrustFin AI is based on a binary classifier interface. Training includes the number of estimators, tree depth, learning rate, row and column subsampling, and a fixed random state.

The practical appeal is the marriage of a tabular classifier and a well-known tree-explanation technique. The drawback is that the decision process that is fitted is distributed over many trees. The motivation for the separate explanation stage comes from reading a single tree, which does not explain the whole ensemble.

\section{Explainable Artificial Intelligence}
\subsection{Local and Global Explanations}
A global explanation describes broader model behaviour, such as which variables tend to be important across a dataset. A local explanation concerns one selected input. These questions are related, but an important global feature need not be the leading reason for every applicant's result.

The implemented TrustFin interface is centred on local review. It receives a contribution vector for the current application and sorts features for display. These bars describe that record under a particular explainer. They should not be presented as an aggregate study of the applicant population or as a discovered rule that always applies to the product.

There is also a distinction between an explanation that is understandable and one that is faithful to the model. A short, familiar reason may be easy to accept while omitting information on which the model depends. Conversely, an accurate model explanation may expose an uncomfortable or undesirable dependence. The literature therefore motivates evaluating explanation behaviour separately from the readability of the presentation.

\subsection{SHAP and Additive Attribution}
Lundberg and Lee introduce SHAP through an additive feature-attribution framework linked to Shapley values \cite{lundberg2017}. A local output is represented through a baseline contribution and feature contributions. The interpretation depends on the model output being explained and on the assumptions used to represent absent features.

TrustFin AI uses TreeExplainer on each XGBoost classifier. The input to this explainer is the transformed feature vector, so the contributions refer to the coordinates that the classifier receives. A category such as property area may appear as several encoded coordinates, while an engineered ratio appears alongside some of the raw quantities used to derive it.

This creates a presentation responsibility. A bar for a standardised financial value is not a raw amount in rupees. A contribution in the binary model's raw margin space is not a change in approval percentage. The report therefore keeps contribution signs and ranking separate from the model's probability output.

\subsection{LIME and Local Approximation}
Ribeiro, Singh and Guestrin propose LIME as a way to approximate a complex classifier locally with an interpretable model \cite{ribeiro2016}. The selected neighbourhood, feature representation, weighting and fitting procedure influence the resulting explanation.

The current wrapper supplies transformed training data to the tabular explainer and passes XGBoost's probability function when explaining a record. Its contributions concern the approval class. A deterministic seed makes repeated requests more repeatable in the same environment, but it does not establish that the sampled neighbourhood is financially realistic.

The distinction is especially relevant for engineered features. If income, principal and FOIR are sampled independently in transformed space, a perturbed point may no longer describe a coherent application. That is a limitation of the implemented explanation setting, even when the original submitted record is valid. It belongs in the interpretation of LIME output rather than being hidden by a repeatable random seed.

\section{Evaluating Explanation Quality}
Yeh et al. explore explanation infidelity and sensitivity in relation to changes in model inputs and outputs \cite{yeh2019}. Alvarez-Melis and Jaakkola investigate explanation robustness and the expectation that similar inputs should have similar explanations \cite{alvarezm2018}. In both cases, the evaluation procedure is important to the meaning of the score.

TrustFin AI employs operational checks tailored to the project. Faithfulness uses top-five local concordance for SHAP and LIME; stability compares SHAP ranks and directions across nearby financial applications; and consistency combines top-five feature overlap with direction agreement. These project adaptations should not be treated as identical to the published measures that motivate them.

One of the important implications is that the same numbers do not necessarily mean the same thing. A rank correlation, a set-overlap ratio and a transformed probability difference may all be displayed between zero and one. Their values remain different answers to different questions. Their combination is a design choice that should be apparent to the reader of the project.

\section{Fairness and Counterfactual Reasoning}
\subsection{Fairness in Supervised Learning}
Fairness requires a defined comparison. Hardt, Price and Srebro define equality of opportunity as the results that eligible individuals in various groups receive \cite{hardt2016}. This is not asking if the probability of the model changes if one applicant's Gender field is changed.

The live TrustFin check is for the second question. It swaps the binary Gender field without changing the other profile. This is a measure of individual sensitivity, not a substitute for a complete fairness assessment.

The current synthetic label-generation rules do not include a Gender-dependent adjustment. Gender remains a model input, so model sensitivity still requires examination. The generated distribution cannot be viewed as a representative picture of real applicants or real discrimination.

\subsection{Counterfactual Reasoning and Decision Support}
Wachter, Mittelstadt and Russell talk about counterfactual explanations as changes that could lead to a different decision \cite{wachter2018}. This is not a description of features that influence the decision, but rather a question of what would have to change if a particular decision system were in place.

TrustFin's affordability calculator is less comprehensive. It adjusts selected financial variables and recalculates instalments and repayment ratios. It does not look for the closest model-approved application or call the classifier for every slider movement. It can support a discussion of repayment capacity, but its output should not be described as a model counterfactual.

The NIST AI Risk Management Framework emphasizes the importance of explanation and fairness in a larger, context-specific risk-management process \cite{nistAirmf2023}. That is the perspective that the project takes when it separates model output, explanation evidence and human review. It has not had an institutional risk-management assessment.

\section{Research Gap and Position of the Project}
Table~\ref{tab:literature-comparison} summarises how the reviewed work is used. The distinction between a source method and the project's adaptation is particularly important for the explanation scores.

\begin{table}[htbp]
\centering\small
\caption{Research foundations and their use in TrustFin AI}\label{tab:literature-comparison}
\begin{tabularx}{\textwidth}{p{3.3cm}XX}
\toprule
Study & Main idea & Use in this project\\\midrule
Hand and Henley \cite{handHenley1997} & Consumer credit classification and its data issues & Context for the loan-assessment problem.\\
Friedman \cite{friedman2001}; Chen and Guestrin \cite{chen2016} & Gradient boosting and efficient tree ensembles & Basis for the selected XGBoost classifiers.\\
Lundberg and Lee \cite{lundberg2017} & Additive feature attribution & Basis for local TreeSHAP explanations.\\
Ribeiro et al. \cite{ribeiro2016} & Locally fitted interpretable approximations & Basis for the tabular LIME wrapper.\\
Yeh et al. \cite{yeh2019}; Alvarez-Melis and Jaakkola \cite{alvarezm2018} & Fidelity, sensitivity and robustness questions & Motivation for checking explanation behaviour; not validation of the custom score.\\
Hardt et al. \cite{hardt2016} & Equality of opportunity & Context for distinguishing group fairness from a Gender-flip test.\\
Wachter et al. \cite{wachter2018} & Counterfactual explanations & Clarifies the limits of the affordability calculator.\\\bottomrule
\end{tabularx}
\end{table}

The project introduces well-known tools to a review process. The connection between product-specific prediction, two local explanation methods, explicit diagnostics and saved assessment records is its contribution. It does not propose a new version of SHAP or LIME, nor does it set up a universal scalar measure of explanation trust.

This role also oversees the evaluation process. Reproduction of the generated labels can be evaluated using classification metrics. Explanation checks can explain behaviour in the tests that have been implemented. Another study would be required to demonstrate that the composite aids reviewers in making better decisions or reliably detecting misleading explanations.

\par\medskip\noindent
This chapter has reviewed the concepts of prediction, explanation and fairness used in TrustFin AI. It has also identified the assumptions behind the implementation and why the project's diagnostics should be understood in the context of their actual procedures.
